% Source: https://artofproblemsolving.com/community/c3h2322728_100th_post_celebration_cp
% Stable identity: aops-topic-2322728-q03
\textbf{The magic squares.} The total is $45$, so each line sums to $15$. The four lines through the centre (middle row, middle column, two diagonals) together contain the centre four times and every other cell once, so $4\cdot15=45+3c$ for the centre $c$, and $c=5$.

The number $1$ cannot be in a corner: a corner lies on a row, a column and a diagonal, and the diagonal through it is $1,5,9$; the row and the column would each need two other numbers adding up to $14$, but apart from $5$ and $9$ the only such pair is $\{6,8\}$. So $1$ is in the middle of an edge, and $9$ is opposite it (on the line through the centre). The edge line containing $1$ has two corners adding up to $14$, so they are $6$ and $8$, in either order. Then everything else is forced (each remaining cell completes a line), and one checks it works, e.g.
\[\begin{array}{ccc}8&1&6\\3&5&7\\4&9&2\end{array}\]
So there are $4\cdot2=8$ magic squares: $4$ positions for the $1$ and $2$ orders of $6,8$. They are the rotations and reflections of the square above.

\textbf{The top middle cell.} It is the middle of an edge, so it holds one of the edge numbers $1,3,7,9$. In the $8$ squares (the $8$ symmetries of the one above), each edge number arrives at the top middle cell under exactly $2$ symmetries. So the top middle cell is $1,3,7,9$ in $2$ squares each, and it is $3$ or $9$ in $4$ of the $8$ squares. The probability of acceptance is $\dfrac48=\boxed{\dfrac12}$.
